\ex{讲义例题\quad 根据电阻支路电流恢复电感电流（第53--56页）}{
如图，电压源 $u(t)$ 为输入，$R_3$ 支路的电流 $y=i_3$ 为输出。取 $x_1=i_1$、$x_2=i_2$，其中 $i_1$ 和 $i_2$ 分别是 $L_1$ 和 $L_2$ 的电流。分析在 $R_1=R_2$、$L_1=L_2$ 时能否由已知的输入和输出唯一确定状态，并讨论参数改变后如何判定。
\begin{center}
\begin{circuitikz}
    \draw (0,0) to[V,l_=$u(t)$] (0,2.75)
        to[R,l^=$R_1$] (1.88,2.75) to[L,l^=$L_1$,i>_=$i_1$] (4.3,2.75)
        -- (6.0,2.75) to[R,l^=$R_3$,i>_=$i_3$] (6.0,0) -- (0,0);
    \draw (4.3,2.75) to[R,l_=$R_2$] (4.3,1.4)
        to[L,l_=$L_2$,i>^=$i_2$] (4.3,0);
    \fill (4.3,2.75) circle[radius=1.2pt] (4.3,0) circle[radius=1.2pt];
    \node[left] at (-.05,1.93) {$+$};
    \node[left] at (-.05,.83) {$-$};
\end{circuitikz}
\end{center}
}

\nt{\textbf{知识点}\quad 节点电流、输出对称性、初态不可辨识、能观矩阵的秩。}

\sol 上方分流节点的电压记为 $v$。由节点电流定律，$y=x_1-x_2$，故 $v=R_3(x_1-x_2)$。两条含电感的回路给出
\[
    L_1\dot x_1=u-R_1x_1-v,\qquad
    L_2\dot x_2=v-R_2x_2 .
\]
当 $R_1=R_2=R,\ L_1=L_2=L$ 时，取 $u=0$、$x_1(t_0)=x_2(t_0)=a$。两支路的状态均为 $a e^{-(R/L)(t-t_0)}$，而 $y(t)$ 恒为零。任取不同的 $a$ 会得到同一条零输出，故状态不完全能观。

一般地，令 $r_i=R_i/L_i$、$q=R_3(1/L_1+1/L_2)$。由以上方程算得
\[
  C=\begin{bmatrix}1&-1\end{bmatrix},\qquad
  CA=\begin{bmatrix}-r_1-q&r_2+q\end{bmatrix},\qquad
  \det\begin{bmatrix}C\\CA\end{bmatrix}=r_2-r_1 .
\]
因此该电路完全能观的条件是 $R_1/L_1\ne R_2/L_2$。只改变一个元件的参数时应重新比较两个比值。

\ex{讲义例题\quad 开关断开后的电容电压能否观测（第57页）}{
下图中，$u(t)$ 是输入电压，右侧端口电压 $y(t)$ 是输出。取 $x_1=u_C$（左正右负）、$x_2=i_L$（向下为正）。开关 $K$ 在 $t_0$ 断开。判断断开后能否仅凭 $u,y$ 恢复两个状态。
\begin{center}
\begin{circuitikz}
    \draw (0,0) to[V,l_=$u(t)$] (0,1.4)
        to[nos,l_=$K$] (0,3) -- (1,3)
        to[R,l^=$R_1$] (3.35,3) -- (5.15,3);
    \draw (1,3) -- (1,2.03) to[C,l_=$C$] (3.35,2.03) -- (3.35,3);
    \node[above] at (2.175,2.24) {$+\quad u_C\quad -$};
    \draw (5.15,3) to[R,l^=$R_2$] (5.15,1.5)
        to[L,l^=$L$,i>_=$i_L$] (5.15,0) -- (0,0);
    \draw (5.15,3) -- (6.65,3) node[ocirc]{};
    \draw (5.15,0) -- (6.65,0) node[ocirc]{};
    \draw[->] (6.65,.25) -- node[right] {$y(t)$} (6.65,2.75);
    \node[left] at (-.05,1.29) {$+$};
    \node[left] at (-.05,.18) {$-$};
    \foreach \P in {(1,3),(3.35,3),(5.15,3),(5.15,0)} {
        \fill \P circle[radius=1.2pt];
    }
\end{circuitikz}
\end{center}
}

\nt{\textbf{知识点}\quad 开关拓扑变化、RC 自由放电、用两组不可区分初态证明能观性不足。}

\sol 断开开关后，$R_1$ 与 $C$ 的左侧公共节点没有外部电流。对该节点列电流方程得
\[
     C\dot x_1+x_1/R_1=0,\qquad
     x_1(t)=x_1(t_0^+)e^{-(t-t_0)/(R_1C)} .
\]
因此两条并联支路给右侧节点的净电流为零。取两组均与断开后理想电路相容的初态，令 $i_L(t_0^+)=0$，但令 $u_C(t_0^+)$ 不同。右侧支路电流及输出电压均可一直为零，电容仍按不同初值放电。由同一输入和相同输出无法反推 $u_C(t_0^+)$，故断开后的电容电压不可观。

\ex{讲义例题\quad 测得一个状态能否恢复另一个？（第58--59页）}{
已知
\[
     \dot x_1=-x_1,\qquad \dot x_2=x_1-2x_2+u,\qquad y=x_1.
\]
输入 $u$ 已知。根据图示模型，判断测量 $y$ 能否唯一确定 $x_1,x_2$。
\begin{center}
\autonomousModeDiagram{$y$}
\end{center}
}

\nt{\textbf{知识点}\quad 可直接测得的状态、未知初态的自由项、输出中缺失的状态方向。}

\sol 因为 $y=x_1$，第一状态随时可读。第二状态为
\[
     x_2(t)=e^{-2(t-t_0)}x_2(t_0)
        +\int_{t_0}^t e^{-2(t-\tau)}[y(\tau)+u(\tau)]\,d\tau .
\]
积分部分由已知信号决定。$x_2(t_0)$ 可以任取而不改变 $y$，所以第二状态无法由 $u,y$ 恢复，系统不完全能观。

\ex{讲义例题\quad 同一衰减速率下的和信号（第60--61页）}{
给定 $\dot x_1=-x_1$、$\dot x_2=-x_2$，只测量 $y=x_1+x_2$。问能否通过任意一段输出记录分别恢复 $x_1(t_0)$ 和 $x_2(t_0)$。
}

\nt{\textbf{知识点}\quad 能观性取决于输出随时间提供的独立信息。}

\sol 对两个标量状态分别积分，有 $y(t)=e^{-(t-t_0)}[x_1(t_0)+x_2(t_0)]$。因此测量记录只给出初态之和。初态 $[a\ \ b]^\top$ 与 $[a+\eta\ \ b-\eta]^\top$ 在每个时刻的输出完全相同。
$CA=-C$，能观矩阵秩为 $1$，不可观方向为 $[1\ \ {-1}]^\top$。
